\documentclass[conference]{IEEEtran}
\IEEEoverridecommandlockouts
\usepackage{cite}
\usepackage{amsmath,amssymb,amsfonts}
\usepackage{booktabs}
\usepackage{array}
\usepackage{graphicx}
\usepackage{textcomp}
\usepackage{xcolor}
\usepackage{url}
\usepackage{balance}
\def\BibTeX{{\rm B\kern-.05em{\sc i\kern-.025em b}\kern-.08em
    T\kern-.1667em\lower.7ex\hbox{E}\kern-.125emX}}

\title{Will Artificial Intelligence Replace Humans?\\
Mapping Automation, Augmentation, and Complementarity Under Uncertainty}

\author{\IEEEauthorblockN{Xcientist Research Team}
\IEEEauthorblockA{Interdisciplinary Artificial Intelligence Research\\
Design-only research report\\
September 2026}}

\begin{document}
\maketitle

\begin{abstract}
The question of whether artificial intelligence will replace humans conflates task automation, occupational change, human augmentation, and the possibility of complementary human-AI teams. This paper proposes a preregistered, design-only benchmark that maps the boundary among these outcomes under routine conditions, distribution shift, and equivocality. The design compares human-only, AI-only, and hybrid conditions using calibrated confidence, selective abstention, and assumption-elicitation interfaces. It measures accuracy, cost-weighted risk, time, calibration, error dependence, delegation, human learning, deskilling, workload, and subgroup robustness. A repeated-exposure component tests whether productivity gains preserve or weaken unaided human performance. A workflow simulation translates task results into task bundles, exception handling, review capacity, and newly created human work without claiming to predict whole occupations. The primary estimands are risk-constrained expected utility per unit time and hybrid complementarity relative to the best individual agent. Evidence from the 2025 Stanford AI Index, a large customer-support study, and task-based labor economics motivates the design. The study does not forecast universal human replacement and reports no empirical result; it provides an auditable protocol for determining where automation, augmentation, or complementarity is scientifically supported.
\end{abstract}

\begin{IEEEkeywords}
artificial intelligence, automation, human-AI collaboration, uncertainty, equivocality, calibration, labor tasks
\end{IEEEkeywords}

\section{Introduction}

Artificial intelligence can perform some computations at speeds beyond human ability, but speed is not the same as reliable judgment. Current systems can classify, summarize, generate code, and optimize decisions at scale; they can also fail under distribution shift, mishandle missing information, express unjustified confidence, or optimize a literal objective that does not capture what a human decision maker means. The question ``Will AI replace humans?'' is therefore scientifically underspecified.

There are at least four distinct claims. A system can automate a narrow task without replacing the human job that contains the task. It can augment a worker by reducing search or documentation time while leaving goal selection and exception handling to the worker. It can form a complementary team when its errors differ from human errors and an interface enables useful delegation. Finally, automation can change the task bundle itself by removing some tasks and creating others in verification, maintenance, and accountability. These claims should not be inferred from one benchmark score.

This paper proposes the \emph{Uncertainty-Calibrated Human-AI Collaboration Benchmark} (U-COLLAB). It turns the broad future question into an experimentally tractable boundary map:

\begin{equation}
\begin{aligned}
&\text{task regime}\times\text{interface}\times\text{expertise}\\
&\qquad\longrightarrow \text{automation, augmentation, or complementarity}.
\end{aligned}
\end{equation}

The benchmark varies routine conditions, distribution shift, and equivocality; compares human-only, AI-only, and hybrid agents; manipulates calibration and abstention; and follows repeated exposure with a no-AI transfer block. It also measures complementary errors and human learning, because a team can be fast but unsafe, or productive today but less capable tomorrow.

The report is explicitly \texttt{DESIGN\_ONLY}. No participants have been enrolled, no model has been evaluated, and no empirical outcome is asserted. The proposed result is a preregistered protocol and interpretation matrix. Its strongest conclusion, if later supported by data, would be conditional: AI automates validated task slices, augments some workers, and complements humans under particular uncertainty and workflow conditions. It would not be a forecast that all humans or all occupations will be replaced.

\section{Background and Evidence}

\subsection{Technical progress is rapid but uneven}

The 2025 Stanford AI Index documents continued gains on demanding benchmarks, falling inference costs, growing organizational adoption, and increasing deployment in medicine, transportation, and other sectors \cite{stanford2025}. It reports that 78\% of organizations surveyed in 2024 used AI, and that inference for a system at roughly GPT-3.5 capability became much cheaper between 2022 and 2024. These trends make task automation economically plausible in more settings.

The same report emphasizes a jagged capability frontier. Systems improved sharply on several demanding benchmarks but continued to struggle with complex reasoning tasks such as structured planning and logic. This is a useful warning against equating average capability with dependable performance in underspecified environments. A system can be very fast on a stable input distribution and still require human supervision when the distribution, objective, or causal structure changes.

\subsection{Evidence for heterogeneous augmentation}

Brynjolfsson, Li, and Raymond studied the staggered introduction of a generative AI conversational assistant among 5,179 customer-support agents \cite{brynjolfsson2025}. The NBER record reports a 14\% average productivity increase measured by issues resolved per hour, a 34\% improvement for novice and low-skilled workers, and minimal impact for experienced and highly skilled workers. It also reports improved customer sentiment, employee retention, and suggestive evidence of worker learning. The published version appears in the *Quarterly Journal of Economics*.

This study is important for two reasons. First, it is evidence that AI can augment people rather than immediately eliminate the workflow. Second, the effect is heterogeneous: novices gain more than experts. A benchmark that reports only an average productivity effect would hide this structure. U-COLLAB therefore treats expertise and repeated exposure as design factors. It asks whether AI transfers useful practice to novices, whether experts remain valuable under shift and ambiguity, and whether repeated assistance changes unaided performance.

The result should not be overgeneralized. Customer support is one workflow, and its task distribution, feedback loop, and language interface may differ from manufacturing, scientific research, management, or care work. The study motivates a measurement strategy, not a universal conclusion.

\subsection{Displacement and new tasks}

Acemoglu and Restrepo formulate automation in terms of the allocation of tasks to capital and labor \cite{acemoglu2019}. Automation can replace labor in existing tasks, producing a displacement effect. New tasks in which labor has a comparative advantage can create a reinstatement effect. Their empirical decomposition associates slower recent employment growth with stronger displacement, weaker reinstatement, and slower productivity growth. Autor similarly argues that automation substitutes for some routine tasks while increasing the relative value of problem formulation, social interaction, and judgment \cite{autor2015}.

These frameworks imply that a job is not a single task and that productivity growth is not identical to labor replacement. A task benchmark should therefore be connected to a workflow model that includes review, exception handling, goal specification, data maintenance, and accountability. U-COLLAB reports task outcomes first and translates them into a simulated task bundle second. This avoids using a task-level win as evidence for occupation-wide substitution.

\subsection{Human-AI symbiosis and automation behavior}

Jarrahi describes human-AI symbiosis in organizational decision making: machines can process large volumes of information while humans contribute contextual interpretation and handling of ambiguity \cite{jarrahi2018}. The scientific version of this claim is complementary error structure. A hybrid is useful when it succeeds on cases that either member would miss, not merely when it averages their outputs.

Human oversight can also fail. Parasuraman and Riley identify automation misuse, disuse, and abuse: people may over-trust an automation, reject useful assistance, or be placed in a workflow that forces unsafe reliance \cite{parasuraman1997}. Confidence displays can help only if they are calibrated and tied to an actionable delegation policy. Amershi \emph{et al.} recommend clear control boundaries, support for correction, and appropriate communication of system limitations \cite{amershi2019}. These principles motivate the interface manipulations in the benchmark.

NIST's AI Risk Management Framework recommends incorporating trustworthiness considerations into the design, development, use, and evaluation of AI systems \cite{nist2023}. U-COLLAB operationalizes that guidance through versioned models, confidence calibration, audit logs, risk constraints, held-out testing, human accountability, and a requirement to restrict deployment when performance under shift is not established.

\section{Scientific Problem and Hypotheses}

The operational question is:

\begin{quote}
Across routine, distribution-shifted, and equivocal task environments, when does AI automation outperform humans, when does human-AI complementarity outperform either alone, and which uncertainty-aware delegation mechanisms prevent over-reliance and loss of human capability?
\end{quote}

We define uncertainty in three parts. Aleatory uncertainty reflects variability in the world that additional data may not remove. Epistemic uncertainty reflects incomplete model knowledge that may be reduced with data or a better model. Equivocality is different: evidence supports multiple plausible interpretations because the objective, meaning, or causal explanation is incomplete. A model can be well calibrated on a stable label distribution while still failing to ask what the task means.

The hypotheses are:

\begin{itemize}
\item \textbf{H0 (performance-only automation):} AI dominates human performance on routine stationary items, but a hybrid does not reliably exceed the best individual agent after time and communication costs are included.
\item \textbf{H1 (distribution-shift boundary):} calibrated confidence and selective abstention reduce cost-weighted loss under shift relative to answer-only AI, especially when confidence tracks error probability.
\item \textbf{H2 (equivocality complementarity):} a hybrid with explicit assumption elicitation outperforms AI-only and human-only conditions on utility and calibration for under-specified tasks.
\item \textbf{H3 (complementary errors):} hybrid value increases when conditional human and AI errors are less correlated, but a persuasive interface can erase this value through automation bias.
\item \textbf{H4 (expertise heterogeneity):} novices obtain larger immediate productivity gains, whereas experts contribute more under shift and equivocality; repeated exposure can improve learning or cause deskilling.
\item \textbf{H5 (uncertainty-aware delegation):} an abstention and clarification policy produces a better risk-coverage curve than fixed always-answer or always-review policies at equal time budget.
\item \textbf{H6 (task-conditional replacement):} strong evidence for replacement is restricted to narrow stationary task slices; workflow-level human tasks persist or are newly created in verification, escalation, objective selection, and accountability.
\end{itemize}

\section{U-COLLAB Benchmark}

\subsection{Factorial structure}

U-COLLAB uses human-only, AI-only, and human-plus-AI conditions. The human study is a proposed repeated-measures design with an illustrative planning target of 240 participants: 80 novices, 80 intermediate participants, and 80 experts. This number is not an observed sample and must be re-estimated by simulation for the smallest preregistered interaction, with attrition included, before data collection. Institutional review, informed consent, privacy protection, and debriefing are required if the study is executed.

\begin{table}[t]
\caption{Primary experimental factors}
\label{tab:factors}
\centering
\scriptsize
\begin{tabular}{p{0.17\columnwidth}p{0.26\columnwidth}p{0.36\columnwidth}}
\toprule
Factor & Levels & Scientific purpose \\
\midrule
Agent & Human, AI, hybrid & Baselines and team value \\
Regime & Routine, shift, equivocal & Locate reliability boundary \\
Interface & Answer, confidence, abstain, assumptions & Test delegation mechanisms \\
Calibration & Calibrated, overconfident, underconfident & Separate truth from persuasion \\
Expertise & Novice, intermediate, expert & Heterogeneous effects \\
Exposure & Single, repeated, no-AI transfer & Learning and deskilling \\
Cost & Symmetric, asymmetric & Risk-sensitive decisions \\
\bottomrule
\end{tabular}
\end{table}

All item generators, answer policies, cost matrices, and interface text are frozen before confirmatory testing. The oracle-calibrated reference supplies a confidence signal tied to held-out error probability. It is a measurement control and is not treated as available deployment technology. The AI checkpoint, prompt or policy, and evaluation domain are versioned.

\subsection{Architecture of the decision interface}

The AI produces a recommendation, a confidence estimate, an abstention option, and a structured list of assumptions. The answer-only interface exposes only the recommendation. The calibrated-confidence interface exposes a probability and a concise evidence summary. The confidence-plus-abstention interface permits the AI to defer. The assumption-elicitation interface asks the human to identify a missing fact or objective that could change the answer before showing the recommendation.

Calibration is manipulated by a preregistered monotone transformation of confidence while holding the answer policy fixed. This separates the value of truthful uncertainty from the persuasive effect of a numerical display. A report-length and time-matched control separates rationale content from extra interaction. Hybrid participants can accept, reject, edit, request clarification, or delegate; they always control final submission.

\subsection{Task families}

Tasks are synthetic, low stakes, and generated from known simulators. They are selected to represent distinct information structures rather than to imitate a single profession.

\begin{enumerate}
\item \textbf{Routine classification and scheduling:} inputs are complete, labels are stable, and the cost matrix is known. The task establishes throughput and basic accuracy.
\item \textbf{Shifted system-state diagnosis:} synthetic equipment or logistics records change feature frequency, missingness, and causal relationships between training and held-out distributions. It tests calibration and transfer without making clinical claims.
\item \textbf{Ambiguous incident interpretation:} a record admits two plausible explanations and the objective is incomplete until a participant asks a fixed clarification question or states an assumption. Unjustified certainty is penalized.
\item \textbf{Asymmetric-cost allocation:} false positives and false negatives have different costs. The task tests whether the agent adapts its threshold to consequences rather than maximizing generic accuracy.
\item \textbf{Counterfactual planning:} the participant predicts how an unchosen intervention changes a simulated system. It tests causal reasoning and whether explanations change decisions for valid reasons.
\end{enumerate}

For equivocal items, the data-generating process contains multiple plausible interpretations. Participants receive credit for identifying ambiguity and acquiring information. A forced label is not assumed to be the only expression of intelligence.

\section{Protocol}

\subsection{Stage S0: preregistration and pilot}

The team preregisters task generators, regime labels, cost matrices, model checkpoint, interface text, calibration transformations, randomization schedule, exclusion rules, primary contrasts, and thresholds for interpretation. A pilot checks comprehension, timing, and technical reliability only; pilot observations are not included in confirmatory estimates.

The initial work is simulation-only. If people later participate, the protocol uses low-stakes decisions, compensation independent of AI acceptance, anonymized logs, and a right to withdraw. No participant is evaluated for employment, credit, health, immigration, or legal status. Professional participants do not provide confidential workplace records.

\subsection{Stage S1: baseline}

Participants complete a human-only block balanced across regimes. They report confidence and can abstain. The block estimates task difficulty, natural calibration, expertise differences, and baseline information acquisition. AI-only episodes use the same item generator but are executed without exposing final answers to participants. Item identities and ground truth remain hidden.

\subsection{Stage S2: collaboration}

Participants receive hybrid blocks in a counterbalanced Latin-square order. Each block combines new item instances with one interface and one calibration condition. The participant may accept, reject, edit, request evidence, or ask for one fixed clarification. The log records recommendation, confidence, assumptions, interaction time, final answer, override, and post-decision confidence.

In equivocal blocks, the key outcome is not just accuracy. It is whether the participant recognizes an under-specified goal, asks for information with high value, and updates the decision when the answer arrives. In shifted blocks, the key outcome is whether confidence and abstention remain aligned with error risk. In routine blocks, the key outcome is whether speed and accuracy justify narrow task automation.

\subsection{Stage S3: repeated exposure and transfer}

The collaboration blocks are repeated on new task instances. The final block removes AI assistance, changes the task distribution, and includes a delayed follow-up where feasible. Learning is an improvement in unaided performance, calibration, and explanation of disagreements relative to baseline. Deskilling is a decline in those outcomes after repeated assistance. Neither result is assumed in advance.

\subsection{Stage S4: workflow bridge}

Item-level results feed a discrete-event workflow simulation. The simulation varies staffing, review capacity, AI latency, error costs, exception rates, and the creation of new human tasks. It includes verification, missing-information collection, objective selection, model monitoring, maintenance, escalation, and accountability. The output is a task-bundle map, not an occupation forecast. Sensitivity analysis reports how conclusions change when review capacity and new-task creation vary.

\section{Measurements and Estimands}

\subsection{Risk-constrained utility}

For agent condition $c$ and regime $r$, define

\begin{equation}
U_{c,r}=\frac{\mathbb{E}[R\mid c,r]}{\mathbb{E}[T\mid c,r]},
\qquad \Pr(L>\tau\mid c,r)\leq\alpha,
\end{equation}

where $R$ is task reward, $T$ is time, and $L$ is cost-weighted loss. The risk constraint prevents a fast but unsafe policy from winning by ignoring rare failures. The primary comparison is the calibrated hybrid against the best individual member within each regime.

Hybrid complementarity is defined as

\begin{equation}
C_r=U_{H+A,r}-\max(U_{H,r},U_{A,r}).
\end{equation}

Positive $C_r$ supports conditional team value. It is not evidence that human-AI teams are universally superior and does not by itself imply that human work should be eliminated.

\subsection{Calibration and selective prediction}

We use Brier score, expected calibration error, reliability diagrams, and risk-coverage curves. If an agent abstains, selective risk at coverage $q$ is

\begin{equation}
SR(q)=\mathbb{E}[L\mid \text{accepted},\ \mathrm{coverage}=q].
\end{equation}

The primary shift comparison tests whether calibrated confidence plus abstention lowers $SR(q)$ at equal coverage and time relative to answer-only AI. In equivocal tasks, we add assumption accuracy and value of information from clarification. A confidence display that is numerically calibrated but ignored is not counted as a successful delegation mechanism.

\subsection{Error dependence and reliance}

For item $i$, let $E_{H,i}$ and $E_{A,i}$ be binary error indicators. We estimate conditional dependence by regime, expertise, and cost structure. The hybrid supplies complementary value when it correctly resolves cases where one member errs. It supplies false redundancy when it accepts correlated wrong answers. We record false acceptance of an incorrect AI recommendation, false rejection of a correct one, time to override, and trust updates after feedback.

\subsection{Human learning and workload}

Unaided transfer measures accuracy, calibration, explanation of disagreement, and retention after delay. Workload combines objective time and interaction count with a short validated subjective scale. The no-AI block is necessary because immediate productivity can increase while the human's independent capability decreases. A result that improves both productivity and transfer is stronger evidence of augmentation than productivity alone.

\subsection{Equity and tail behavior}

Difficulty, missingness, and cost asymmetry vary across synthetic task strata. If future studies collect demographic data, this requires explicit consent and a separate fairness analysis. Report subgroup effects by expertise and relevant task strata rather than making speculative claims about social groups. A mean improvement is insufficient if tail loss rises for a vulnerable task slice.

\section{Statistical Analysis}

The confirmatory model is hierarchical with fixed effects for agent, regime, interface, calibration, expertise, exposure, cost, and preregistered interactions. Random intercepts and slopes cover participant, task family, item generator, and session. A generic outcome model is

\begin{equation}
\begin{aligned}
Y_{pijk}={}&\beta_0+\beta_C C_i+\beta_R R_j+\beta_I I_k\\
&+\beta_{CR}C_iR_j+\beta_{CI}C_iI_k+u_p+u_f+u_g+\epsilon_{pijk}.
\end{aligned}
\end{equation}

Primary contrasts are AI versus human on routine items, calibrated hybrid versus AI on shifted items, assumption-elicitation hybrid versus both baselines on equivocal items, and repeated-exposure transfer with versus without AI. Correct for the preregistered primary-outcome family. Bootstrap at the participant and task-family level; repeated decisions from one participant are not independent observations.

Before recruitment, simulate the full design under plausible effect sizes and intraclass correlations. The illustrative target of 240 is accepted only if the smallest expertise-by-regime-by-interface interaction has adequate power after attrition. If recruitment or effect assumptions change, the amendment is recorded before unblinding confirmatory outcomes.

The boundary decision is conditional. **Automation** is supported for a narrow task slice when AI meets the risk constraint, remains calibrated on held-out data, and does not create unacceptable shift loss. **Augmentation** is supported when hybrid utility or human learning improves without a harmful increase in reliance. **Complementarity** is supported when $C_r$ is positive, error dependence is sufficiently low, and the effect survives expertise and interface sensitivity checks. None of these states implies occupation-wide replacement.

\section{Threats to Validity}

The first threat is benchmark overfitting. Hidden generators, changed surface forms, new distributions, and an independent replication reduce it. The second is interface confounding: a rationale may consume time or persuade independently of truth. Time-matched text controls and answer-only conditions estimate that effect. The third is automation bias: a confidence number may become an authority signal. Miscalibrated-confidence and abstention controls measure this failure rather than assuming it away.

The fourth threat is expertise selection. Volunteers may be unusually favorable or skeptical about AI. Stratified recruitment, preregistered targets, and a follow-up sample from a different population provide robustness. The fifth is learning contamination: participants may memorize item answers. New generators and changed surface forms reduce memorization. The sixth is ecological validity. Synthetic tasks provide causal control but cannot reproduce every organizational incentive, social relationship, or liability structure. The workflow bridge is therefore a sensitivity tool, not proof of labor-market impact.

The seventh threat is model drift. A model can change after deployment, or an interface can cause users to explore a distribution not represented in the test. Versioned checkpoints, domain declarations, periodic recalibration, and an abstention default are required. The eighth is selective reporting. All task-family failures, missing trials, safety exclusions, and human overrides are included. A result that survives only after removing difficult cases is not a boundary map.

\section{Governance and Responsible Deployment}

\section{Boundary Map and Workflow Translation}

\subsection{From task scores to task allocation}

The benchmark's output is a boundary map, not a leaderboard. For each task regime, the team reports the best risk-constrained individual, the hybrid gain, the review cost, and the change in unaided human capability. A task is a candidate for narrow automation only when the AI's expected utility satisfies the loss constraint on held-out items and when domain drift is monitored. A task is an augmentation candidate when the human-plus-AI condition improves utility or learning without increasing harmful reliance. A task is complementary when the hybrid resolves cases that are conditionally difficult for either member.

Let $d_r$ denote the proportion of task instances delegated to AI in regime $r$, and let $v_r$ be the value of human review after the AI's confidence and assumptions are observed. A selective delegation policy chooses

\begin{equation}
d_r^*=\arg\max_{d\in[0,1]}\{U_r(d)-\lambda L_r(d)-\gamma T_r(d)\},
\end{equation}

where $U_r$ is workflow utility, $L_r$ is cost-weighted loss, $T_r$ is coordination time, and $\lambda,\gamma$ are preregistered policy weights. The policy is not optimized after observing the test set. Its purpose is to show how the same AI can be useful at one coverage level and harmful at another.

\subsection{Routine regime}

Routine items have stable semantics, complete information, and a known cost matrix. If AI performance exceeds human performance with lower time and the risk constraint is met, a narrow automation claim is defensible. Even here, monitoring and exception handling remain tasks. The relevant question is not whether a person presses a button for each item, but who maintains the data pipeline, checks model drift, updates the cost matrix, and handles cases outside the validated distribution.

\subsection{Distribution-shift regime}

Shifted items change feature frequencies, missingness, or causal relationships. A model may preserve a high average score while becoming overconfident on a rare but costly slice. The benchmark evaluates calibration, selective risk, and abstention at fixed coverage. A hybrid gain is meaningful if human review is directed toward cases for which the AI is uncertain or systematically wrong, rather than simply asking the human to recheck every answer. If the AI cannot identify its own shift, the workflow must reduce delegation coverage or pause.

\subsection{Equivocal regime}

Equivocal items contain multiple plausible interpretations or incomplete objectives. They are not merely noisy classification problems. The participant can state an assumption or request one fixed clarification. A system that produces a confident label without identifying the missing objective may be fast but not intelligent in the relevant operational sense. The benchmark credits information acquisition and objective clarification, then measures whether the final choice changes for the correct reason. This is a direct test of the limitation in the topic prompt without treating human intuition as an unmeasurable essence.

\subsection{Task-bundle simulation}

The workflow bridge converts item-level results into a directed task graph. Nodes represent input acquisition, model inference, human review, clarification, escalation, decision, monitoring, maintenance, and accountability. Edges carry arrival rates, latency, and error probabilities estimated from the benchmark. A discrete-event simulation varies review capacity and the rate at which new exception tasks appear. The output includes throughput, tail loss, human time, AI coverage, and the share of tasks that remain human-controlled.

This translation exposes two opposite errors in public discussion. First, high AI accuracy on one node does not prove that the workflow can be automated because upstream data and downstream accountability may dominate. Second, a hybrid can reduce total labor on a validated slice while creating more valuable verification or objective-setting work. Those newly created tasks are part of technological change, not evidence that the human has been replaced in every sense.

\subsection{Decision boundary visualization}

The intended figure can be reconstructed from released tables by plotting risk-constrained utility against AI coverage for each regime. The automation region is the portion in which AI-only utility meets the loss bound. The augmentation region is where hybrid utility or learning exceeds the best individual without a reliance penalty. The complementarity region is where $C_r$ is positive and conditional error dependence is low. Regions are reported separately for expertise strata and cost asymmetry; a single global boundary would conceal precisely the heterogeneity the study is designed to reveal.

\section{Robustness and Reproducibility}

\subsection{Stress tests before deployment}

The evaluation includes four stress-test families. Sensor and feature missingness test whether the system communicates epistemic uncertainty. Temporal drift changes the frequency of events and the cost of mistakes. Objective perturbations change which outcome is preferred while keeping the observations similar. Interface removal tests whether the human can still make a safe decision when confidence or rationale is unavailable. Each stress test is held out from training and labeled by the type of uncertainty it introduces.

The red team also tests persuasive failure. It presents a wrong recommendation with high displayed confidence, a correct recommendation with low confidence, and a plausible but unsupported rationale. The outcome is not merely acceptance rate; it is the relationship between acceptance, confidence calibration, and final loss. A safe system should make it possible to reject a persuasive wrong answer and should not punish a human for asking for clarification.

\subsection{Model and interface versioning}

The AI checkpoint, prompt or policy, tokenizer, retrieval data, calibration transform, interface build, and task generator receive version identifiers. The audit record links every decision to those identifiers. A future model update is a new experimental condition, not a silent continuation of the old result. Calibration is re-estimated on a time-separated validation set, and deployment coverage is reduced whenever the measured error distribution changes.

The release package includes raw item definitions, hidden labels under controlled access, response logs, code for the utility and calibration metrics, randomization seeds, and the workflow simulator. Confirmatory scripts read a frozen manifest and fail if an unregistered task or outcome is added. Exclusions, missing trials, participant withdrawals, and technical failures are reported in the denominator. These controls prevent a favorable boundary map from being produced by selective deletion.

\subsection{Independent replication}

An independent team receives the frozen model interface and task specification but generates a new item bank. It evaluates the primary contrasts, error dependence, and no-AI transfer without access to the original participant-level results. Replication is required for a strong complementarity claim because interface details and local training can produce spurious team gains. A result that survives a second implementation, new task generator, and altered ordering is more credible, but it remains conditional on the tested workflow.

\subsection{Sensitivity to cost and review capacity}

Utility depends on consequences. The paper therefore reports a grid of symmetric and asymmetric cost matrices and varies review capacity in the workflow simulation. If the recommended delegation policy changes under a modest cost perturbation, the result is a sensitive boundary and should not be summarized as a robust replacement claim. If the policy is stable across cost and review ranges, the automation or augmentation interpretation is stronger. The analysis never hides a rare severe error inside an average score.

\subsection{Human capability as a state variable}

A workflow model usually treats human skill as fixed, but AI assistance can change it. U-COLLAB includes a state variable for unaided human calibration and transfer. Repeated exposure updates that state according to observed practice, feedback, and time spent on exceptions. A policy that delegates all easy items may improve current throughput but reduce opportunities for human practice. Conversely, a policy that assigns explanation and correction tasks may preserve learning while still reducing search time. This dynamic is essential for deciding whether a short-run augmentation effect is sustainable.

\subsection{Interpretation checklist}

Before using the result to justify a deployment or a claim about replacement, investigators must answer six questions: Was the task distribution declared? Did the AI meet the risk constraint under held-out shift? Was confidence calibrated and acted upon? Were human and AI errors conditionally complementary? Did unaided human ability improve or decline? Which tasks remained or were created in the workflow? A missing answer means the evidence supports at most a narrower task claim.

The protocol preserves human accountability for hybrid decisions. Audit logs contain model version, data provenance, confidence, abstention, clarification requests, overrides, final decisions, and detected shift. The system displays its validated domain and does not present an AI recommendation as an authority. If shift loss, tail risk, or calibration failure exceeds a preregistered limit, the workflow defaults to human review or pauses deployment.

The design forbids unreviewed high-stakes use. No medical, legal, employment, credit, immigration, or safety-critical decision is made in an initial study. Human-subject data are minimized, anonymized, and retained only as needed for the preregistered analysis. Researchers disclose the model, limitations, cost matrix, and likely failure modes. NIST AI RMF provides a governance vocabulary for validity, reliability, transparency, explainability, privacy, fairness, and accountability \cite{nist2023}.

The benchmark also treats the loss of human capability as a system risk. If a workflow obtains a short-term productivity gain while human workers lose the ability to detect errors, the system is not classified as successful augmentation. Training, no-AI practice, review rotation, and escalation authority become part of the deployment design. The organization must decide which objectives remain human-controlled, especially when the objective is equivocal rather than merely uncertain.

\section{Expected Contributions and Limitations}

\section{Expected Result Patterns}

Because this paper is a design-only protocol, the following are prospective interpretation patterns rather than results. Their purpose is to prevent a future analysis from treating every AI advantage as evidence for replacement.

\subsection{Narrow automation}

A narrow automation pattern would occur if AI meets the risk constraint, remains calibrated on held-out routine items, and achieves higher utility per unit time with less review than human-only performance. This pattern justifies a limited task-scope claim. The validated data distribution, cost matrix, and coverage must be specified. It does not justify deployment on shifted or equivocal items, and it does not show that a job bundle has disappeared. Monitoring, maintenance, exception handling, and accountability remain part of the workflow.

\subsection{Conditional augmentation}

An augmentation pattern would occur if the hybrid improves utility or throughput while preserving or improving unaided human calibration and transfer. The result is stronger when novices receive useful practice and experts retain authority over exceptions and goals. If immediate productivity increases while no-AI transfer declines, the result is a mixed pattern: short-run augmentation with a possible long-run deskilling cost. The deployment response would be training, practice rotation, and explicit review rather than simply increasing AI coverage.

\subsection{Complementary team performance}

A complementarity pattern requires more than a higher hybrid average. The hybrid must resolve cases that are conditionally hard for either member, with low or explainable human-AI error dependence. It must survive changes in interface, expertise, cost asymmetry, and item generator. If a hybrid succeeds only because humans check every AI output, it may be a costly review system rather than a scalable complement. If it succeeds because calibrated abstention routes uncertain cases to the right reviewer, the delegation policy itself is the scientifically important result.

\subsection{Uncertainty failure}

An uncertainty failure pattern would occur if AI remains fast but becomes overconfident under shift, or if it confidently chooses one interpretation when the objective is equivocal. A confidence number that does not predict errors is not a reliable control signal. A rationale that persuades participants to accept wrong answers is a risk factor, even if the rationale reads fluently. In this pattern, the workflow should reduce AI coverage, require clarification or human review, and preserve an audit trail. The result would reveal a boundary of safe use, not a general failure of artificial intelligence.

\subsection{Labor reallocation}

The workflow simulation may show that automation of one node reduces the time spent on routine classification while increasing demand for verification, escalation, objective selection, monitoring, and maintenance. This is a task reallocation result. Its magnitude depends on review capacity, arrival rates, model latency, and the creation of new tasks. The simulation should therefore report sensitivity bands instead of one point forecast. An outcome in which the original task is automated but the workflow still requires human judgment is neither a simple replacement nor a claim that AI has no economic effect.

\subsection{What would falsify the proposal}

The proposed mechanism would be weakened if calibration interfaces do not change delegation despite truthful confidence; if assumption elicitation does not help on equivocal items; if hybrid performance is no better than the best individual after time and risk costs; if human and AI errors are strongly correlated across all regimes; or if all apparent learning gains disappear in independent no-AI transfer. The study would also fail as a boundary map if results depend on a single model version, visible test labels, or post hoc task exclusions. These are falsification conditions for the benchmark's claims, not reasons to reinterpret an unfavorable outcome as hidden intelligence.

\subsection{Reporting standard}

Every future empirical paper using this design should report the full task-regime vector, risk constraint, coverage, calibration curves, error dependence, human overrides, no-AI transfer, subgroup effects, and workflow assumptions. It should state which claims are task-level, workflow-level, or economic extrapolations. This reporting standard makes it possible to compare a new model with an older one without confusing capability progress with safe replacement. It also makes negative evidence valuable: a reproducible failure under equivocality can guide interface and training research.

U-COLLAB makes five contributions. First, it separates task automation from job replacement. Second, it distinguishes distribution shift from equivocality, where clarification and objective selection matter. Third, it defines complementarity through risk-constrained team utility and conditional error dependence. Fourth, it measures the temporal effect of assistance on learning, deskilling, and over-reliance. Fifth, it links results to a workflow task-bundle simulation and a governance protocol.

The possible outcomes are all informative. If AI wins only on routine tasks, the result identifies a narrow automation slice. If calibrated abstention improves shifted performance, uncertainty communication is a real intervention rather than a cosmetic feature. If assumption elicitation produces hybrid gains on equivocal cases, human value is partly located in objective and context handling. If repeated assistance helps novices but harms no-AI transfer, augmentation has a learning trade-off. If human and AI errors are correlated, deploying both may create false redundancy.

The protocol cannot forecast all occupations or all future systems. Human intuition is operationalized as calibrated uncertainty handling, information acquisition, and goal clarification; this may omit tacit knowledge that is difficult to encode. Calibration is distribution-specific. A good interface may fail under different incentives, liability rules, or privacy constraints. The proposed sample and task generators are planning choices, not empirical facts.

The largest conceptual limit is that replacement is not a scalar property of AI. It is an emergent outcome of technology, task design, institutions, skill formation, wages, and newly created work. The benchmark can identify functional and workflow boundaries under specified assumptions. It cannot prove that humans as a class will or will not be replaced.

\section{Conclusion}

Artificial intelligence will likely automate some tasks and augment some workers, but the evidence does not support a universal binary forecast. The scientifically useful question is where, under what uncertainty regime, and with which delegation interface AI changes the allocation of tasks and the performance of human-AI teams. U-COLLAB addresses that question with matched human-only, AI-only, and hybrid conditions; calibration and abstention manipulations; equivocality tasks; error-dependence analysis; repeated exposure; and a workflow bridge.

The responsible conclusion is conditional. AI speed can support narrow automation when risk and distribution are controlled. Human-AI complementarity is plausible when errors differ and humans retain authority to question objectives and assumptions. In uncertain or equivocal environments, calibrated abstention and clarification may be more valuable than an always-answer system. A future empirical run could refine these boundaries, but this design-only report makes no observed claim and does not predict universal human replacement.

\balance
\begin{thebibliography}{00}
\bibitem{stanford2025} Stanford Institute for Human-Centered Artificial Intelligence, \emph{AI Index Report 2025}, Stanford University, 2025.
\bibitem{brynjolfsson2025} E. Brynjolfsson, D. Li, and L. R. Raymond, ``Generative AI at work,'' \emph{Quarterly Journal of Economics}, vol. 140, no. 2, pp. 889--942, 2025, doi: 10.1093/qje/qjae044.
\bibitem{acemoglu2019} D. Acemoglu and P. Restrepo, ``Automation and new tasks: How technology displaces and reinstates labor,'' \emph{Journal of Economic Perspectives}, vol. 33, no. 2, pp. 3--30, 2019, doi: 10.1257/jep.33.2.3.
\bibitem{autor2015} D. Autor, ``Why are there still so many jobs? The history and future of workplace automation,'' \emph{Journal of Economic Perspectives}, vol. 29, no. 3, pp. 3--30, 2015, doi: 10.1257/jep.29.3.3.
\bibitem{jarrahi2018} M. H. Jarrahi, ``Artificial intelligence and the future of work: Human-AI symbiosis in organizational decision making,'' \emph{Business Horizons}, vol. 61, no. 4, pp. 577--586, 2018, doi: 10.1016/j.bushor.2018.03.007.
\bibitem{parasuraman1997} R. Parasuraman and V. Riley, ``Humans and automation: Use, misuse, disuse, abuse,'' \emph{Human Factors}, vol. 39, no. 2, pp. 230--253, 1997, doi: 10.1518/001872097778543886.
\bibitem{amershi2019} E. Amershi \emph{et al.}, ``Guidelines for human-AI interaction,'' in \emph{Proceedings of the CHI Conference}, 2019, doi: 10.1145/3290605.3300233.
\bibitem{nist2023} National Institute of Standards and Technology, \emph{Artificial Intelligence Risk Management Framework (AI RMF 1.0)}, NIST AI 100-1, 2023, doi: 10.6028/NIST.AI.100-1.
\end{thebibliography}

\end{document}
